\documentclass[10pt,a4paper]{amsart}
\usepackage[margin=19mm]{geometry}
\usepackage{amsmath,amssymb,mathtools}
\usepackage[scr=rsfs,cal=euler]{mathalfa}
\usepackage{xfrac}
\usepackage[unicode,hidelinks,bookmarks=false]{hyperref}
\newcommand{\ie}{i.e.\ }
\newcommand{\cf}{cf.\ }
\newcommand{\resp}{respectively\ }
\newcommand{\Subsection}{\subsection}
\providecommand{\nobreakdash}{\nobreak\hbox{-}}
\setlength{\emergencystretch}{2em}
\multlinegap0pt
\usepackage{mathtools}
\usepackage{booktabs}
\usepackage{tikz}
\usepackage{cleveref}
\newtheorem{theorem}{Theorem}[section]
\newtheorem{lemma}[theorem]{Lemma}
\newtheorem{corollary}[theorem]{Corollary}
\newcommand{\Deck}{\mathcal D}
\newcommand{\Aut}{\operatorname{Aut}}
\newcommand{\Even}{\operatorname{Even}}
\begin{document}
\title[Nonisomorphic Graphs Can Share an Arbitrarily Large Fraction]{Nonisomorphic Graphs Can Share an Arbitrarily Large Fraction of Their Vertex-Deleted Cards}
\author{Sergey Ivanov}
\date{}
\maketitle
\noindent\emph{Source and licence.} Nonisomorphic Graphs Can Share an Arbitrarily Large Fraction of Their Vertex-Deleted Cards. Version arXiv:2608.11930v1 (CC BY 4.0). This material is distributed under the Creative Commons Attribution 4.0 International licence, http://creativecommons.org/licenses/by/4.0/, as recorded in the arXiv OAI licence metadata for this exact version and shown on the abstract page https://arxiv.org/abs/2608.11930v1. Full rights notice: licenses/arxiv-2608.11930-CC-BY-4.0.txt.\par\smallskip
\noindent\textit{Editorial scope.} Counted unit: the original \texttt{theorem} environment \texttt{thm:main} at source lines 290--305 together with its complete proof at lines 306--326; the delivered document keeps source lines 209--327 so that every referenced label and every citation resolves inside it. Native editable source is the paper's own concatenated TeX; the AMS wrapper is editorial formatting only. No publisher class, upstream script or unrelated artwork is redistributed.
\section{General construction and proof}\label{sec:general}
The four-port example already contains the complete mechanism. To make the common-card fraction closer
to $1$, we use more port classes while retaining a selector that distinguishes even from odd
permutations. Fix an even integer $r=2k\ge4$.

We first need a seed whose own symmetries are all even. On labeled vertices
$A=\{a_1,\ldots,a_k\}$ and $D=\{d_1,\ldots,d_k\}$, let $X_r$ consist of a clique on $A$, an independent
set $D$ (no two vertices of $D$ are adjacent), and the private edges $a_i d_i$. A symmetry can
arbitrarily permute the $k$ matched pairs
$(a_i,d_i)$, but it must apply the same permutation once on $A$ and once on $D$. Its sign is therefore
squared, so it is even on all $r$ labels. In particular, $X_r$ has $k!$ symmetries and all lie in
$\Even(r)$.

Let
\[
 \mathcal F_r=\{\sigma X_r:\sigma\in\Even(r)\}.
\]
Construct a selector $S_r$ with ports $p_i$, pair vertices $q_{ij}$, and a clique of vertices $z_F$
indexed by $F\in\mathcal F_r$. As before, port $p_i$ is adjacent to the pair vertices containing $i$,
and $z_F$ is adjacent to $q_{ij}$ exactly when $ij\in E(F)$. Since $X_r$ has $k!=(r/2)!$
symmetries,
\[
 |\mathcal F_r|=\frac{r!}{2(r/2)!},\qquad
 C_r=\binom r2+\frac{r!}{2(r/2)!}
\]
is the number of nonport selector vertices.

The next lemma is the general version of the four-port parity check. Its proof has two parts: degrees
identify the three layers, and incidence then forces every port permutation to preserve the chosen
family $\mathcal F_r$.

\begin{lemma}[Permutations allowed by the general selector]\label{lem:selector}
The port permutations induced by automorphisms of $S_r$ are precisely the even permutations of $[r]$.
\end{lemma}
\begin{proof}
First, the permutations preserving $\mathcal F_r$ are exactly the even ones. Indeed, every even
permutation preserves the family by definition. Conversely, if a permutation $\tau$ preserves
$\mathcal F_r$, then $\tau X_r=\alpha X_r$ for some even $\alpha$. Thus
$\alpha^{-1}\tau$ is a symmetry of $X_r$ and is even by the preceding construction, so $\tau$ is even.

It remains to recognize the layers inside $S_r$. Ports have degree $r-1$. Every pair vertex occurs in
the same number $\lambda$ of graphs in $\mathcal F_r$, so its degree is $2+\lambda$, while a selector
vertex has degree $|\mathcal F_r|-1+|E(X_r)|$. For $r=4$ these degrees are $3,5,8$. For $r=2k\ge6$,
\[
 \lambda=|\mathcal F_r|\frac{k+1}{2(2k-1)},
\]
which gives $2+\lambda>r-1$; also $\lambda<|\mathcal F_r|$ and $|E(X_r)|\ge3$, so selector vertices
have still larger degree. An automorphism therefore preserves all three layers. Port incidence forces
$q_{ij}\mapsto q_{\sigma(i)\sigma(j)}$, and selector incidence forces $F\mapsto\sigma F$. Hence
$\sigma$ preserves $\mathcal F_r$ and is even. Conversely, every even permutation extends to all three
layers.
\end{proof}

We now store numerical information in the ports. For a vector
$\mathbf a=(a_1,\ldots,a_r)$ with every $a_i\ge2$, let $S_r(\mathbf a)$ replace $p_i$ by an independent
false-twin class $P_i$ of size $a_i$. Because these classes are structurally recognizable, an
isomorphism must match their sizes using a permutation allowed by the selector.

\begin{lemma}[Recognizing the blown-up ports]\label{lem:blowup}
$S_r(\mathbf a)\cong S_r(\mathbf b)$ if and only if an even permutation $\sigma$ satisfies
$a_i=b_{\sigma(i)}$ for every $i$.
\end{lemma}
\begin{proof}
The $P_i$ are the only nontrivial false-twin classes: distinct classes have distinct pair-vertex
neighborhoods, pair vertices are identified by their incident classes, and selector vertices have
distinct pair-vertex neighborhoods. Any isomorphism therefore permutes the $P_i$. Contracting them
recovers $S_r$, so \cref{lem:selector} makes the induced permutation even. Conversely, an allowed even
permutation extends by arbitrary bijections between corresponding classes.
\end{proof}

The final size pattern uses consecutive integers and swaps only the first two entries. This makes the
intact assignment odd, but a deletion from almost any later class creates a repeated size and hence a
second, harmless transposition.

For $t\ge1$, define
\[
 \mathbf a(t)=(t+1,t+2,\ldots,t+r),\qquad
 \mathbf b(t)=(t+2,t+1,t+3,\ldots,t+r),
\]
and set $G_{r,t}=S_r(\mathbf a(t))$ and $H_{r,t}=S_r(\mathbf b(t))$.


\begin{theorem}[Fixed-$r$ overlap bound]\label{thm:main}
For every even $r\ge4$ and $t\ge1$, the graphs $G_{r,t},H_{r,t}$ are connected, nonisomorphic, have the
same degree sequence, and have
\[
 n_{r,t}=rt+\frac{r(r+1)}2+C_r
\]
vertices. Moreover,
\[
 b(G_{r,t},H_{r,t})\ge L_{r,t}
   =(r-1)t+\frac{r(r+1)}2-1,
\]
and therefore
\[
 \liminf_{t\to\infty}\frac{b(G_{r,t},H_{r,t})}{n_{r,t}}\ge1-\frac1r.
\]
\end{theorem}

\begin{proof}
The unique permutation matching the two intact size vectors is the odd transposition
$\tau=(1\;2)$, so \cref{lem:blowup} gives $G_{r,t}\not\cong H_{r,t}$.

For $j=2$, delete from $P_2$ in $G_{r,t}$ and from $P_1$ in $H_{r,t}$; the two resulting vectors agree
and yield $t+2$ common cards. For $j\in\{3,\ldots,r\}$, delete from $P_j$ on both sides. Positions
$j-1$ and $j$ in the first card then have equal size, so $\rho_j=(j-1\;j)$ is harmless. The product
$\tau\rho_j$ is even and, by \cref{lem:blowup}, gives an isomorphism of the cards. Thus
\[
 b(G_{r,t},H_{r,t})\ge\sum_{j=2}^r(t+j)
 =(r-1)t+\frac{r(r+1)}2-1.
\]
The port classes contain $rt+r(r+1)/2$ vertices and the selector contributes $C_r$, proving the order.
Connectivity is immediate. Port degrees and selector degrees agree in both graphs; pair-vertex degrees
are $a_i+a_j+\lambda$, whose multiset is unchanged when the size vector is permuted. Hence the degree
sequences agree. Finally, with $r$ fixed,
\[
 \frac{L_{r,t}}{n_{r,t}}
 =\frac{(r-1)t+O_r(1)}{rt+O_r(1)}\longrightarrow1-\frac1r.
\]
\end{proof}


\end{document}
